\documentclass[10pt,oneside]{book}

% Page geometry - 125 × 200 mm (Gyldendal 1999 Tractatus format).
\usepackage[
  paperwidth=125mm,
  paperheight=200mm,
  top=20mm,
  bottom=18mm,
  left=28mm,
  right=12mm,
  headheight=12pt,
  headsep=6mm,
  footskip=8mm,
]{geometry}

\usepackage{fontspec}
\usepackage{polyglossia}
\setdefaultlanguage[variant=nynorsk]{norwegian}

\IfFontExistsTF{EB Garamond}{
  \setmainfont{EB Garamond}[
    Numbers={Lining,Proportional},
    SmallCapsFeatures={LetterSpace=4},
  ]
}{
  \setmainfont{TeX Gyre Pagella}
}
\renewcommand{\normalsize}{\fontsize{10.5}{12.5}\selectfont}
\normalsize

\usepackage{microtype}

\newfontfamily{\mathfbfont}{Cambria Math}[Scale=0.95]
\newcommand{\msym}[1]{{\mathfbfont #1}}

\usepackage{titlesec}
\titleformat{\chapter}[block]
  {\normalfont\centering\scshape\addfontfeature{LetterSpace=10}\large}
  {}{0em}{\MakeLowercase}
\titlespacing*{\chapter}{0pt}{2.5em}{2.5em}

\titleformat{\section}[block]
  {\normalfont\centering\scshape\addfontfeature{LetterSpace=6}\normalsize}
  {}{0em}{\MakeLowercase}
\titlespacing*{\section}{0pt}{2.5em}{1.5em}

\titleformat{\subsection}[block]
  {\normalfont\scshape\addfontfeature{LetterSpace=2}\small}
  {}{0em}{\MakeLowercase}
\titlespacing*{\subsection}{0pt}{1.5em}{0.8em}

\renewcommand{\thechapter}{}

\usepackage{fancyhdr}
\pagestyle{fancy}
\fancyhf{}
\fancyhead[L]{\small\scshape\leftmark}
\fancyhead[R]{\small\thepage}
\fancyheadoffset[L]{16mm}
\renewcommand{\headrulewidth}{0.4pt}
\renewcommand{\footrulewidth}{0pt}

\fancypagestyle{plain}{
  \fancyhf{}
  \fancyhead[L]{\small\scshape\leftmark}
  \fancyhead[R]{\small\thepage}
  \fancyheadoffset[L]{16mm}
  \renewcommand{\headrulewidth}{0.4pt}
  \renewcommand{\footrulewidth}{0pt}
}

\fancypagestyle{titlepage}{
  \fancyhf{}
  \fancyheadoffset[L]{0pt}
  \renewcommand{\headrulewidth}{0pt}
  \renewcommand{\footrulewidth}{0pt}
}

\newcommand{\silentchapter}[2]{%
  \clearpage
  \markboth{#2}{}%
  \addcontentsline{toc}{chapter}{#1}%
}

\makeatletter
\renewcommand\tableofcontents{%
  \markboth{innhald}{}%
  \begingroup\footnotesize
  \@starttoc{toc}%
  \endgroup
}
\renewcommand*\l@chapter[2]{%
  \ifnum \c@tocdepth >\m@ne
    \vskip 0pt%
    \setlength\@tempdima{1.5em}%
    \begingroup
      \footnotesize
      \parindent \z@ \rightskip \@pnumwidth
      \parfillskip -\@pnumwidth
      \leavevmode \normalfont
      \advance\leftskip\@tempdima
      \hskip -\leftskip
      #1\nobreak\leaders\hbox{$\m@th\mkern \@dotsep mu\hbox{.}\mkern \@dotsep mu$}\hfill
      \nobreak\hb@xt@\@pnumwidth{\hss #2}\par
    \endgroup
  \fi}
\makeatother

\usepackage{calc}
\newlength{\propnumwidth}
\setlength{\propnumwidth}{13mm}
\newlength{\propnumgap}
\setlength{\propnumgap}{2mm}

\newcommand{\prop}[3]{%
  \par\addvspace{6pt}%
  \noindent\llap{\makebox[\propnumwidth][l]{\textbf{#1}\textsuperscript{\textit{#2}}}\hspace{\propnumgap}}%
  #3\par%
  \addvspace{2pt}%
}

\newcommand{\propcont}[1]{%
  \par\addvspace{3pt}%
  \noindent #1\par%
}

\newenvironment{widebody}{%
  \begin{adjustwidth}{-16mm}{0mm}%
}{%
  \end{adjustwidth}%
}

\usepackage{changepage}
\usepackage{graphicx}
\graphicspath{{analysis/figures/}}
\usepackage{tikz}

\begin{document}
\frontmatter

\thispagestyle{titlepage}
\begin{tikzpicture}[remember picture, overlay]
  \node[anchor=north west, inner sep=0] at ([xshift=12mm, yshift=-18mm]current page.north west) {%
    {\fontsize{36}{40}\selectfont\bfseries FORMLÆRE}%
  };
  \node[anchor=north west, inner sep=0] at ([xshift=12mm, yshift=-32mm]current page.north west) {%
    {\large\itshape Ei traktatform om korleis form oppstår}%
  };
\end{tikzpicture}
\clearpage

\silentchapter{Føreord}{føreord}
\begin{widebody}
\begingroup\small\setlength{\parskip}{4pt}
Om kvifor ting har den forma dei har. Svaret som har dominert den designteoretiske diskursen gjennom store delar av det tjuande hundreåret, er formulert i Louis Sullivans doktrine: form følgjer funksjon. Denne påstanden er ikkje berre upresis; han er fundamentalt villfartande. Analysar av den materielle historia har demonstrert at form aldri følgjer funksjon på ein deterministisk måte; form følgjer form.\footnote{Michl (1995)} Kvar ny konfigurasjon er ein modifikasjon av ein eksisterande konfigurasjon, framkalla av agentar som arbeider innanfor strengt avgrensa moglegheiter og med inkomplett kunnskap om heilskapen. Funksjonalismen hindra den teoretiske forståinga av korleis formgjeving faktisk fungerer, ettersom han postulerte at det finst éin optimal løysingsveg for eit gjeve problem; eit ekko av ein platonsk ideologi.

George Kubler kom nærare ein adekvat morfologi. Han synte korleis former ordnar seg i sekvensar, der kvar konfigurasjon utgjer eit punkt i eit formelt forløp.\footnote{Kubler (1962)} Eit fåtal punkt fungerer som primærobjekt som opnar mutasjonar og nye sekvensar, medan den overveldande majoriteten er replikasjonar. Kubler handsama forma si historie som eit problem definert av posisjon og tid, strippa for den reine funksjonalismens teleologi. Likevel mangla han eit formelt apparat for å logisk forklare dei underliggjande kreftene som tvingar desse sekvensane i spesifikke retningar. Han skildra den observerte rørsla, men han identifiserte ikkje vektorane som genererte ho.

Denne traktaten etablerer eit substrat-uavhengig rammeverk for korleis form oppstår. Teksten forklarer ikkje isolert sett kvifor ein spesifikk gjenstand har fått sin spesifikke fasong, men utleier den generative prosessen som systematisk produserer alle former. Den grunnleggjande tesen er konstaterande: Form er ein posisjon i eit multidimensjonalt rom av moglege konfigurasjonar, forma av samtidige seleksjonstrykk og navigert av agentar på fleire skalaer.

Traktaten skildrar posisjonar og krefter, ikkje estetisk verdi. Det er mogleg å forklare logisk kvifor to ulike konfigurasjonar okkuperer ulike posisjonar i same landskap, men den estetiske dommen ligg strengt utanfor grensene for dette systemet. Modellen registrerer utelukkande den resulterande posisjonen og vektorane som førte dit; intensjonen og den indre kognitive opplevinga til formgjevaren er irrelevante for den formelle analysen av topologien.

Denne traktaten lukkast berre i den grad ho let seg fellast. Proposisjonane i denne teksten tener som eit logisk stillas; når formgjevaren har erkjent mekanikken i landskapet, kan stillaset kastast, og intuisjonen kan atter ta over, men no forankra i eksakt kunnskap om kvifor intuisjonen i det heile fungerer.
\endgroup
\end{widebody}

\silentchapter{Ordliste}{ordliste}
\begin{widebody}
\begingroup\small
\ordlisteentry{Affordanse}{Dei formene og handlingane eit materiale tilbyr og motset seg. Avgjer kva formoperasjonar som er fysisk moglege i eit gjeve substrat.}
\ordlisteentry{Agensspektrum}{Ein skala som indikerer at agentar varierer i korleis dei best let seg styre: frå direkte maskinvaremodifikasjon til utveksling av abstrakte grunnar.}
\ordlisteentry{Agent}{Eit system uttømande definert av trippelet: mål, måling og justering. Agens krev utelukkande ei funksjonell organisering basert på negativ tilbakekopling.}
\ordlisteentry{Formalgebra}{Mengda av former innanfor eit euklidsk rom, utstyrt med operasjonane sum, differanse og produkt. Grunnlaget for formgrammatikken.}
\ordlisteentry{Formgjevar}{Kvar instans av agens i formdomenet. Marknaden, materialet og handverkaren er alle formgjevarar.}
\ordlisteentry{Formgrammatikk}{Ein algebra over former i eit euklidsk rom, bygd på reglar som opererer under vilkår av innleiring.}
\ordlisteentry{Formrom (Morphospace)}{Mengda av alle moglege konfigurasjonar innanfor ein gjeven klasse.}
\ordlisteentry{Innleiring (Embedding)}{Ein euklidsk transformasjon som beviser at éi form eksisterer som ein integrert del av ei anna form.}
\ordlisteentry{Kanalisering}{Graden av robustheit til ein konfigurasjon mot forstyrringar. Målt som brattleiken på dalveggane i landskapet.}
\ordlisteentry{Klasse}{Ei gruppering av objekt som deler den eksakt same hovudfunksjonen.}
\ordlisteentry{Kognitiv lyskjegle}{Den romleg-temporale regionen ein agent evnar å representere og påverke.}
\ordlisteentry{Konfigurasjon}{Ein bestemt samanheng av objekt. Forma er konfigurasjonens struktur.}
\ordlisteentry{Objekt}{Den enklaste bestanddelen i ein form. Objektet er udeleleg og uforanderleg.}
\ordlisteentry{Projeksjon}{Den matematiske avbildinga av ei form inn i eit formrom. Iboande tapsbringande.}
\ordlisteentry{Seleksjonstrykk}{Ein faktor som systematisk endrar sannsynlegheitsfordelinga over formrommet.}
\ordlisteentry{Substrat}{Det fysiske eller algoritmiske mediet som realiserer agensen.}
\ordlisteentry{Tilpassingslandskap}{Den aggregerte verknaden av alle samtidige seleksjonstrykk over formrommet.}
\endgroup
\end{widebody}

\silentchapter{Innhald}{innhald}
\begin{widebody}
\tableofcontents
\end{widebody}

\mainmatter

\addcontentsline{toc}{chapter}{1 Formverda er alt som er tilfelle}
\markboth{formverda er alt som er tilfelle}{}
\prop{1}{}{Formverda er alt som er tilfelle.}
\prop{1.1}{o}{Formverda er totaliteten av realiserte konfigurasjonar, ikkje av ting.}
\prop{1.11}{t}{Alt som har form, har nettopp denne forma og ikkje ei anna.}
\prop{1.12}{t}{At ein konfigurasjon tek ein eksakt form, når uendeleg mange andre var logisk moglege, krev ei forklaring utover konfigurasjonen sjølv.}
\prop{1.13}{t}{Forklaringa ligg i relasjonen mellom den manifesterte konfigurasjonen og det totale settet av konfigurasjonar som var moglege, men ikkje vart realiserte.}
\prop{1.2}{d}{Eit objekt er den enklaste bestanddelen i ein form. Objektet er udeleleg og uforanderleg. Det utgjer formverdas substans.}
\prop{1.21}{d}{Ein konfigurasjon er ein bestemt samanheng av objekt. Forma er konfigurasjonens struktur.}
\prop{1.22}{t}{Sidan objektas natur inneber moglegheita for å inngå i konfigurasjonar, er alle moglege former allereie gjeve i og med objekta.}
\prop{1.3}{d}{Formrommet (morphospace) til ein klasse er mengda av alle moglege konfigurasjonar for objekt i denne klassen.}
\prop{1.31}{d}{Kvart realisert objekt utgjer eitt eksakt punkt i dette n-dimensjonale rommet.}
\prop{1.32}{d}{Det empiriske formrommet er alltid ein projeksjon. Inga endeleg mengd parametrar captures den latente kompleksiteten fullt ut.}
\prop{1.4}{d}{Formrommet deler seg topologisk i tre regionar: dei busette, dei opne og dei forbodne.}
\prop{1.41}{o}{Dei busette regionane utgjer det historiske arkivet. Dei fungerer som ankerpunkt for all framtidig navigasjon og utgjer den induktive premissen for vidare form.}
\prop{1.42}{o}{Dei opne regionane utgjer det tilstøytande moglege. Dei er teoretisk tilgjengelege, men enno ikkje aktualiserte.}
\prop{1.43}{t}{Dei forbodne regionane er stengde av fysiske og materielle lovar.}
\prop{1.44}{t}{Grensa for det forbodne er ein funksjon av teknologien ved eit gjeve historisk augeblink.}
\prop{1.5}{o}{Konfigurasjonane i formrommet er ikkje uniformt fordelte. Dei samlar seg i klynger med store tomrom mellom seg.}

\clearpage
\addcontentsline{toc}{chapter}{2 Trykket er det som endrar sannsynet}
\markboth{trykket er det som endrar sannsynet}{}
\prop{2}{}{Trykket er det som endrar sannsynet.}
\prop{2.1}{t}{Av 1.5: Asymmetrien i formrommet provar at det finst gradientar. Desse gjer visse konfigurasjonar meir sannsynlege enn andre.}
\prop{2.11}{t}{Gradientar er historisk betinga.}
\prop{2.2}{t}{Sannsynlegheitsfordelinga for formgenerering er isomorf med ein ergodisk Markoff-prosess.}
\prop{2.21}{t}{Sannsynlegheita for ei ny form er betinga av dei umiddelbart føregåande tilstandane i nettverket.}
\prop{2.3}{t}{Sjølvforsterkande mekanismar skaper sti-avhengigheit. Eit tidleg, tilfeldig val etablerer ein standard.}
\prop{2.31}{t}{Symmetribrot er uunngåeleg. Små historiske hendingar tippar dominansen mot éi av to likeverdige formelle løysingar.}
\prop{2.4}{t}{Kvar realisert form er eit poly-agentisk kompromiss. Sidan fleire motstridande krav verkar samstundes, er inga form optimal for alle.}
\prop{2.5}{d}{Materialet legg sine eigne sannsynsfordelingar over rommet gjennom affordansar.}
\prop{2.51}{d}{Affordansar er operasjonar materialet logisk og strukturelt tillét eller motset seg.}
\prop{2.6}{o}{Nye substrat arvar formspråket til det materialet dei erstattar, heilt til agentens kunnskapsrom har integrert dei nye affordansane.}

\clearpage
\addcontentsline{toc}{chapter}{3 Landskapet er den aggregerte topologien}
\markboth{landskapet er den aggregerte topologien}{}
\prop{3}{}{Landskapet er den aggregerte topologien.}
\prop{3.1}{d}{Den aggregerte verknaden av alle kognitive og materielle krav skaper eit tilpassingslandskap.}
\prop{3.11}{d}{Kvar konfigurasjon har ei høgd i landskapet som svarar til kor godt forma tilfredsstiller gradientane.}
\prop{3.2}{t}{Landskapet er ruglete. Det er karakterisert av fleire lokale maksima.}
\prop{3.21}{t}{På eit ruglete landskap vert adaptive vandringar fanga av lokale haugar.}
\prop{3.22}{t}{Å forlate ein lokal haug krev at agenten fyrst aksepterer ei dårlegare tilpassa form.}
\prop{3.3}{t}{Trykka er hierarkiske. Dei sterkaste trykka graver djupe kanalar. Proporsjonar her er robuste mot endring.}
\prop{3.4}{d}{Ein stil er ei klynge av konfigurasjonar samla kring ein lokal attraktor.}
\prop{3.41}{t}{Stilar har uskarpe grenser fordi landskapet er eit topologisk kontinuum.}
\prop{3.5}{t}{Seleksjon er eit filter, ikkje ei skapande kraft. Ho eliminerer konfigurasjonar som fall for djupt i dalane.}
\prop{3.6}{t}{Konvergent formgjeving oppstår når uavhengige tradisjonar navigerer langs same gradient mot same fjelltopp.}

\clearpage
\addcontentsline{toc}{chapter}{4 Tida er landskapets bevegelse}
\markboth{tida er landskapets bevegelse}{}
\prop{4}{}{Tida er landskapets bevegelse.}
\prop{4.1}{a}{Tilpassingslandskapet deformerer seg i takt med teknologi, marknad og samfunn.}
\prop{4.11}{t}{Når landskapet vert deforma, kan stabile haugar brått verte djupe dalar. Agentar må vandra, elles døyr dei ut.}
\prop{4.2}{t}{Endringsraten er karakterisert av avbrote likevekt. Morfologisk stase er normaltilstanden.}
\prop{4.21}{t}{Stase vert avbroten av episodar med ekstrem radiasjon inn i nyopna regionar.}
\prop{4.3}{t}{Det moglege vert falda med det realiserte. Kvar aktualisert form opnar eit nytt tilstøytande mogleg rom.}
\prop{4.4}{o}{Materialstraumar er den kraftigaste mekanismen for landskapsdeformasjon.}
\prop{4.5}{t}{Ko-evolusjon krev balanse på randa av kaos. Dersom landskapet deformerer seg for raskt, kollapsar systemet.}

\clearpage
\addcontentsline{toc}{chapter}{5 Agens er evna til å navigere}
\markboth{agens er evna til å navigere}{}
\prop{5}{}{Agens er evna til å navigere.}
\prop{5.1}{d}{Formgjeving vert driven av agentar. Ein agent er definert ved tilbakekoplingssyklusar: evna til å kjenne avstand til eit mål og kapasitet til å justere.}
\prop{5.2}{d}{Kvar agent navigerer med ei avgrensa kognitiv lyskjegle.}
\prop{5.21}{d}{Lyskjegla definerer den regionen av formrommet som agenten kan representere og utføre operasjonar innanfor.}
\prop{5.3}{d}{Navigasjonen er den formelle operasjonen mellom Konseptrommet (C) og Kunnskapsrommet (K).}
\prop{5.31}{d}{Eit konsept er ein uavgjord proposisjon; han er korkje sann eller usann i eksisterande kunnskap. Utan uavgjord status finst berre deduksjon, ikkje design.}
\prop{5.32}{t}{Design er den gjensidige utvidinga av K og C.}
\prop{5.4}{t}{Agentar treng ikkje innsikt i det totale landskapet. Dei opererer på den lokale gradienten.}
\prop{5.5}{d}{Læring er den temporale utvidinga av agentens lyskjegle.}
\prop{5.6}{d}{Navigasjonen er operasjonell. Mengda av former i eit rom utgjer ein formalgebra.}
\prop{5.61}{d}{Den fundamentale relasjonen i formalgebraen er innleiring. Å sjå er å identifisere ei innleiring.}
\prop{5.62}{d}{Ein formgrammatikk er ein algebra der reglane opererer under innleiring.}

\clearpage
\addcontentsline{toc}{chapter}{6 Forma oppstår mellom agentane}
\markboth{forma oppstår mellom agentane}{}
\prop{6}{}{Forma oppstår mellom agentane.}
\prop{6.1}{t}{Ingen einskild agent har absolutt diktat over forma. Forma materialiserer seg der alle seleksjonstrykk balanserar kvarandre ut.}
\prop{6.2}{t}{Materialet er ein passiv, men absolutt agent. Dialogen med materialets indre struktur er premissen for form.}
\prop{6.3}{o}{Agentar utfører nisjekonstruksjon. Dei endrar ikkje berre seg sjølve, men dei romlege vilkåra for alle omliggjande agentar.}
\prop{6.4}{t}{Den akkumulerte effekten av nisjekonstruksjon utgjer ein ekologisk arv av seleksjonstrykk.}
\prop{6.5}{t}{Resultatet av ein designprosess er per definisjon ikkje dekka av startvilkåra.}

\clearpage
\addcontentsline{toc}{chapter}{7 Sanninga er tekstens evne til å falle}
\markboth{sanninga er tekstens evne til å falle}{}
\prop{7}{}{Sanninga er tekstens evne til å falle.}
\prop{7.1}{t}{Sidan landskapet kontinuerleg vert deforma, kan form aldri nå eit absolutt endemål.}
\prop{7.2}{t}{Ei realisert form er eit fryst augeblink i ein sti-avhengig prosess. Forma er ein mellombels kvilestad.}
\prop{7.3}{t}{Det finale kvilepunktet eksisterer ikkje utanfor termodynamisk likevekt.}
\prop{7.8}{o}{Denne traktaten er sjølv ein konfigurasjon i eit intellektuelt formrom. Stigen vert ikkje kasta; ho vert kalibrert. At teksten kan fellast, er garantien for hennar gyldigheit.}

\appendix
\silentchapter{Referansar}{referansar}
\begin{widebody}
\begingroup\small
\refentry{Sullivan, L. H. (1896). The tall office building artistically considered. Lippincott's Magazine, 57, 403–409.}
\refentry{Thompson, D'A. W. (1917). On growth and form. Cambridge University Press.}
\refentry{Wiener, N. (1948). Cybernetics. MIT Press.}
\refentry{Waddington, C. H. (1957). The strategy of the genes. George Allen \& Unwin.}
\refentry{Kubler, G. (1962). The shape of time. Yale University Press.}
\refentry{Gibson, J. J. (1979). The ecological approach to visual perception. Houghton Mifflin.}
\refentry{Stiny, G. (2006). Shape: Talking about seeing and doing. MIT Press.}
\refentry{Hatchuel, A. \& Weil, B. (2009). C-K design theory. Research in Engineering Design.}
\refentry{Fields, C. \& Levin, M. (2022). Competency in navigating arbitrary spaces. Entropy.}
\endgroup
\end{widebody}

\end{document}
